\section{Evaluasi}
\label{sec:evaluasi-bab9}

\subsection{Soal Pilihan Ganda}

\begin{enumerate}
    \item Pernyataan yang benar tentang direktif assembler adalah
          \dots
          \begin{enumerate}
              \item direktif menghasilkan opcode seperti instruksi
              \item direktif hanya berlaku saat perakitan dan tidak
                    menghasilkan kode mesin
              \item direktif dieksekusi prosesor saat program
                    berjalan
              \item direktif tidak berpengaruh pada berkas
                    \file{.exe}
          \end{enumerate}

    \item Direktif yang menyediakan dua byte per nilai adalah \dots
          \begin{enumerate}
              \item \dir{DB} \item \dir{DW}
              \item \dir{DD} \item \dir{DT}
          \end{enumerate}

    \item Deklarasi \lstasm|nilai DW 1234h| menempatkan byte di
          memori dengan urutan \dots
          \begin{enumerate}
              \item \hex{12} lalu \hex{34}
              \item \hex{34} lalu \hex{12}
              \item \hex{00} lalu \hex{12}
              \item \hex{43} lalu \hex{21}
          \end{enumerate}

    \item Pasangan direktif yang benar beserta sifatnya adalah
          \dots
          \begin{enumerate}
              \item \dir{EQU} dapat didefinisikan ulang, \code{=}
                    tidak dapat
              \item \dir{EQU} tidak dapat didefinisikan ulang,
                    \code{=} dapat
              \item keduanya tidak dapat didefinisikan ulang
              \item keduanya selalu dapat didefinisikan ulang
          \end{enumerate}

    \item Tugas yang dikerjakan linker, bukan assembler, adalah
          \dots
          \begin{enumerate}
              \item menghitung \eng{offset} label di dalam satu
                    segmen
              \item menerjemahkan mnemonic menjadi opcode
              \item menyelesaikan simbol eksternal dan melakukan
                    \eng{relocation}
              \item membaca direktif \dir{EQU}
          \end{enumerate}
\end{enumerate}

\subsection{Soal Uraian}

\begin{enumerate}
    \item Jelaskan perbedaan instruksi mesin dan direktif assembler,
          lalu berikan dua contoh direktif yang mempengaruhi isi
          berkas \file{.exe} dan dua contoh yang tidak.

    \item Tuliskan deklarasi data untuk: string \code{'PRAKTIKUM'},
          \eng{array} berisi empat nilai 32-bit \hex{00000001}, dan
          penampung \eng{byte} berisi enam belas nilai nol. Hitung
          total byte dan \eng{offset} setiap label.

    \item Jelaskan alur translasi \file{.asm} $\rightarrow$
          \file{.obj} $\rightarrow$ \file{.exe}. Sebutkan apa yang
          sudah diketahui assembler pada akhir lintasan pertama dan
          apa yang baru ditentukan linker. Kaitkan jawaban Anda
          dengan penanda \code{R} pada berkas \file{.lst}.

    \item Sebuah program dirakit tanpa pesan kesalahan, tetapi
          linker menolak berkas objeknya. Sebutkan tiga
          kemungkinan penyebabnya beserta cara memperbaikinya.
\end{enumerate}

\subsection{Kisi-kisi Penilaian}

\begin{table}[H]
    \centering
    \caption{Kisi-kisi evaluasi Bab~\ref{chap:translasi}.}
    \label{tab:kisi-bab9}
    \small
    \begin{tabular}{@{}p{1.5cm}p{5.0cm}p{2.6cm}c@{}}
        \toprule
        \textbf{Sub-CPMK} & \textbf{Indikator yang diukur} &
        \textbf{Bentuk soal} & \textbf{Bobot} \\
        \midrule
        2.3 & Membedakan instruksi dan direktif serta
        memilih direktif data & Pilihan ganda 1, 2 & 15\% \\
        \addlinespace
        2.3 & Menghitung ukuran dan \eng{offset} deklarasi data &
        Uraian 2 & 25\% \\
        \addlinespace
        2.3 & Menerapkan aturan \eng{little-endian} dan perbedaan
        \dir{EQU} dengan \code{=} & Pilihan ganda 3, 4; Uraian 1 &
        25\% \\
        \addlinespace
        2.3 & Menjelaskan tugas assembler dan linker serta membaca
        berkas hasil & Pilihan ganda 5; Uraian 3, 4 & 35\% \\
        \bottomrule
    \end{tabular}
\end{table}

\begin{checklistbox}
\textbf{Checklist Pencapaian Kompetensi Bab~\ref{chap:translasi}}

Seorang mahasiswa dinyatakan kompeten pada Sub-CPMK 2.3 apabila:
\begin{itemize}
    \item[\checked] dapat menjelaskan mengapa direktif tidak
          menghasilkan kode mesin;
    \item[\checked] dapat menuliskan deklarasi \dir{DB}, \dir{DW},
          \dir{DD}, dan \dir{DUP} beserta perhitungan ukurannya;
    \item[\checked] dapat menentukan urutan byte
          \eng{little-endian} pada data multi-byte;
    \item[\checked] dapat menjelaskan perbedaan \dir{EQU} dan
          \code{=} serta akibat definisi ulang;
    \item[\checked] dapat menjelaskan tugas assembler dan linker
          pada alur \file{.asm} $\rightarrow$ \file{.obj}
          $\rightarrow$ \file{.exe};
    \item[\checked] berhasil merakit dan mengaitkan program dengan
          TASM dan TLINK, serta membaca berkas \file{.lst} dan
          \file{.map}.
\end{itemize}
\end{checklistbox}
